\subsection{六方密堆积\ch{Co}的铁磁态态密度和能带计算}
在\textrm{VASP}自旋极化计算中，设置计算控制参数\texttt{ISPIN}=2，电子将按照指定的极化方向(一般默认为$z$轴方向)分为自旋向上和自旋向下两部分，对应的\textrm{Kohn-Sham}方程将在各自不同的空间里求解，因此两种自旋的本征态可以有不同的本征值。由于全部格点都有相同的极化方向，因此也称为共线磁结构\textrm{(collinear magnetic structure)}。

不难看出，共线磁结构的极化方向原则上是任意的，本质上只须考虑两个格点的磁交换作用，所以也只能区分铁磁有序和反铁磁有序。如果考虑\textrm{SOC}或非共线磁结构\textrm{(noncollinear magnetic structure)}，则允许每个格点上的极化方向不在同一方向，计算时需要执行\textrm{vasp\_ncl}。

对于六方密堆积\ch{Co}，结构文件\textrm{POSCAR}

根据六方晶格对称性，\textrm{POSCAR}中含有两个\ch{Co}原子，控制参数\texttt{direct}表示原子位置采用分数坐标\textrm{(fractional coordinates)}。

\textrm{INCAR}

\textrm{INCAR}中，\textrm{SCF}计算控制参数
\begin{itemize}
	\item \texttt{ISPIN}=2
	\item \texttt{MAGMOM}中指定\ch{Co}的初始磁矩为$2\mu_{\mathrm{B}}$
	\item 为后续重启计算需要，设置控制参数\texttt{LMAXMIX}，对\textit{d}轨道设置为4，对\textit{f}轨道设置为6；否则写入\textrm{CHGCAR}文件的电荷密度是不完整的，而且当\texttt{LMAXMIX}指定轨道角动量后，\textrm{DFT}迭代计算中的电荷密度混合中，会包含这些轨道，因此加速收敛过程。
	\item 对于磁性计算，有必要设置\texttt{LASPH}，它会在计算过程中对\textrm{PAW}方法在原子附近的计算引入非球形贡献的梯度校正。
	\item \texttt{ISMEAR}=2，属于取到\textrm{Methfessel-Parton}方法展开级数的2阶，对于金属体系这是较好的积分逼近方案。
\end{itemize}
\textrm{DOS}计算的输入文件
\begin{itemize}
	\item \texttt{ICHARG}=11,电荷密度从\textrm{CHGCAR}直接读入，后续计算中电荷密度将保持不变。在自旋极化计算中，\textrm{CHGCAR}中的电荷密度是总的电荷密度。
	\item \texttt{ISMEAR}=-5，对于\textrm{DOS}计算来说，带\textrm{Bl\"ochl}校正的四面体积分得到的\textrm{Fermi}面更精确，而\textrm{DOS}轮廓线也会更平滑，不过收敛得会慢一些。
	\item \textcolor{red}{\texttt{NEDOS}确定了计算\textrm{DOS}的格点数目}
	\item \texttt{LORBIT}=11，对于\textrm{DOS}来说，除了总\textrm{DOS}，在原子附近，按照角动量$l$和磁角动量$m$投影的分波\textrm{DOS}数据，也将写入\textrm{DOSCAR}文件，而分波磁矩也将写入\textrm{OUTCAR}文件中。
\end{itemize}

\textrm{KPOINTS}
在\textrm{SCF}和\textrm{DOS}计算中，$\vec k$空间网格都是均匀剖分的。根据六方密堆积的晶格结构\textrm{POSCAR}文件，可以看到
\begin{itemize}
	\item 晶格矢量$\mathbf{a}_1$和$\mathbf{a}_2$，满足$|\mathbf{a}_1|=|\mathbf{a}_2|$;
	\item 晶格矢量$\mathbf{a}_1$的方向为$(1,0,0)$，$\mathbf{a}_2$的方向为$(-1/2,\sqrt3/2,0)$;
	\item 晶格矢量$\mathbf{a}_3$的方向为$(0,0,1)$，但其晶格长度满足$|\mathbf{a}_3|/|\mathbf{a}_1|\approx1.6$
	\item \textcolor{red}{(利用\textrm{POSCAR}的\texttt{scale}参数的作用，同一个\textrm{POSCAR}文件有两种等价的形式。)}
\end{itemize}
应用\textrm{py4vasp}查看六方密堆积\ch{Co}的初基原胞和原子。

根据实空间晶格和倒空间晶格的对应关系，实空间的晶格越长，倒空间晶格越短。因此在$\vec k$空间网格剖分时，应该保证在$x,y$方向与$z$方向的剖分比保持近似$1.6:1$，使得倒空间网格分布均匀。

绘制能带用的$\vec k$空间网格点是沿高对称性路径选取的。

\textrm{SCF}计算过程的总能和磁矩计算结果可以在\textrm{OSZICAR}中看到。总的磁矩与实验值$1.7\mu_{\mathrm{B}}$/\ch{Co}接近。

因为初基原胞中包含两个\ch{Co}原子，不难看出，总的磁矩已经被平均分配到两个\ch{Co}原子上，因此计算值虽然比实验值略小，但还是吻合得比较好。

虽然\textrm{SCF}计算中的$\vec k$空间积分采用的是\textrm{Methfessel-Paxton}方法并得到比较可靠的结果，但是计算\textrm{DOS}，还是建议使用四面体方法，而如果要计算分波\textrm{DOS}，控制参数\texttt{LORBIT}$\geqslant10$，而电荷密度要保持不变。

采用\textrm{py4vasp}绘制自旋极化的\textrm{DOS}和$d$电子的分波\textrm{DOS}。

从\textrm{DOS}和分波\textrm{DOS}看，体系的磁矩主要来自于\ch{Co}的$d$轨道，自旋向上和自旋向下部分的\textrm{DOS}总体相似，但由于两者相对\textrm{Fermi}能级发生移动，因此最终不同自旋的电子数也有差别。\footnote{在中文口语习惯中，不同自旋的电子态，常常人为地成为``自旋向上''，``自旋向下''加以区分；在英文文献中，习惯用``多数自旋''\textrm{(majority-spin)}和``少数自旋''\textrm{(minority spin)}加以区分。}在经典的\textrm{Stoner}磁学模型中，将这两种不同自旋的能量相对移动，称为``交换裂分''\textrm{(exchange splitting)}。

\textcolor{red}{采用\textrm{py4vasp}模块，可以进一步分析计算结果和相关数据。}

\textcolor{red}{关于局域磁矩的讨论}
在\textrm{OUTCAR}文件中，搜索关键词\textcolor{red}{\textrm{magnetization}}，直到最后，可以看到有关局域磁矩的信息

与\textrm{py4vasp}提取的信息对比

在每个\ch{Co}上，都有关于$s,p,d$轨道的局域自旋磁矩。总的磁矩来自各个原子的局域自旋磁矩的总和。由于\textrm{VASP}仅计算了自旋磁矩的贡献，这里并未包括轨道的轨道磁矩的贡献。另一方面，$d$轨道的轨道磁矩比起自旋磁矩小很多，而对$f$轨道，轨道磁矩与自旋磁矩的贡献就会相当。\textcolor{red}{这将在后续介绍。}

\textcolor{red}{原子的在位磁矩与体系磁矩的关系}
\begin{displaymath}
	m_{\mathrm{atom}=gm\mu_{\mathrm{B}}
\end{displaymath}
这里$g$是$g$因子，$m$是原子的磁量子数。$\mu_{\mathrm{B}}$是\textrm{Bohr}磁子。该磁矩是通过中子散射测得的磁矩。而根据电子自旋计算得到的有效磁矩为
\begin{displaymath}
	m_{\amthrm{atom}}^{\mathrm{eff}}=g\sqrt{s(s+1)}\mu_{\mathrm{B}}
\end{displaymath}
其中$s$是自旋量子数。有效磁矩可以通过$\mu$子散射实验或按照\textrm{Curie}定理，测量比热直接得到。

可以通过\textrm{py4vasp}模块实现局域磁矩的可视化。

如果不考虑SOC，原子的自旋磁矩间可以认为无相互作用。因此计算过程中并不涉及电子磁矩方向的信息，换言之，每个原子自旋结构的可以统一在$z$方向。\textcolor{red}{对于六方密堆积初级原胞中两个\ch{Co}原子来说，由参数\texttt{MAGMOM}设定原子的初始磁矩后，最后得到的分波局域磁矩数值符号相同，因此确定体系的在位磁矩间是铁磁耦合。}

最后计算六方密堆积体系的能带结构。

并通过\textrm{py4vasp}模块查看计算的能带。

可以看到\textrm{DOS}计算中类似的\textrm{Stoner}模型中的交换裂分。
与\textrm{DOS}中类似，与自旋向上的电子本征值比，自旋向下的能量相对于\textrm{Fermi}能向上移动。如能带图中所示，在\textbf{M}点和\textbf{K}点附近，自旋向上的能量有三个能带的本征值在$-0.5\mathrm{eV}$附近，而自旋向下的电子能带中，相应部分在$1.3~\mathrm{eV}$附近。这就是交换裂分。

\subsection{反铁磁的体心立方\ch{Cr}的磁性原胞}




